
\subsection{Instance normalisation as an ensemble-diversity device}

One further ablation is worth reporting because its sign is determined by a
property of the fleet rather than by tuning. Adding a fourth input channel,
the window differenced from its own engine baseline and divided by that
engine's standard deviation, leaves single-model accuracy essentially
unchanged (FD001 $-0.50\%$, FD003 $+0.15\%$, FD002 $+1.46\%$ on official
test) but decorrelates the errors. Five such members added to the fifteen-model
ensemble lower FD001 from 9.64 to 9.47 and FD003 from 10.68 to 10.60.

The same channel is catastrophic on FD004 (validation RMSE 14.85 against 9.5
for the baseline) and does not help FD002. The reason is mechanical: on fleets
with six operating conditions the per-engine variance is dominated by condition
switching, which is an order of magnitude larger than the degradation drift, so
dividing by it compresses precisely the signal the model needs. On
single-condition fleets the per-engine variance is mostly inter-unit nuisance,
and removing it helps. The decision rule is therefore conditioned on the number
of operating conditions, which is known before any test data is seen.
